% !TEX root = ../main.tex
\setlength{\absparsep}{18pt}
\begin{resumo}
Este texto de qualificação examina a dinâmica de uso e cobertura da terra em
Goiás entre 1985 e 2024 e sua relação com fatores socioeconômicos, a partir da
série anual do MapBiomas (Coleção 10.1) para os 246 municípios goianos e de uma
malha de 166 Áreas Mínimas Comparáveis que garante território constante nas
análises longitudinais. Uma periodização orientada pelos dados divide os 40
anos em três atos (pastagem como herança, 1985--2000; expansão e
intensificação, 2001--2019; e conversão acelerada sob rótulo ambíguo,
2020--2024) e organiza o achado central: uma reorganização espacial da
conversão de terras em direção ao norte do estado (``a marcha ao norte''). A
investigação decompõe-se em quatro frentes: a existência do padrão
(deslocamento do centro de massa de cada uso da terra, que sobrevive a
verificações de pixel, malha, classe e barra de erro, com um gradiente
latitudinal que persiste nos quarenta anos); o mecanismo local, em que duas
lógicas convivem em todo o estado e a região explica pouco da idade da pastagem
convertida, ainda que a transição dominante seja de pastagem para lavoura no
Sul e de vegetação natural para pastagem no Norte; o teste formal do
deslocamento causal intraestadual, que não se confirma e tem como leitura
alternativa uma resposta coordenada a um \emph{driver} macroeconômico comum,
distribuída ao longo do gradiente Sul--Norte (evidência corroborante, não
estabelecida); e a desaceleração recente, compatível com o esgotamento da terra
ainda exposta à conversão no Sul, sustentada pela eliminação de hipóteses
concorrentes e por um mecanismo medido, a queda da taxa de conversão à medida
que a região se esgota. Apresentam-se a metodologia, os resultados
preliminares, os limites do que o trabalho afirma e o cronograma até a defesa.

\textbf{Palavras-chave}: uso e cobertura da terra; Cerrado; fronteira
agrícola; Goiás; MapBiomas.
\end{resumo}
